% !PRINTED 222
\weberpage{222}

%<<CONT
dessen Wurzeln die $x_{h,0}$ sind. Der Einfluß einer Substitution der Gruppe
$\mathfrak{A}_{1,0}$
\[
\begin{pmatrix} a, & b \\ 0, & \partial \end{pmatrix}
\]
auf $x_{h,0}$ besteht darin, daß $x_{h,0}$ in $x_{\partial h,0}$ übergeht, und die
Gruppe dieser Gleichung $\varphi(n)$ten Grades besteht daher aus den Vertauschungen
\[
(x_{h,0},\ x_{\partial h,0}),
\]
worin $\partial$ jede beliebige zu $n$ teilerfremde Zahl sein kann. Diese Gruppe ist
eine \gesperrt{Abelsche} und daher sind die Wurzeln $x_{h,0}$ \gesperrt{nach Adjunktion von $\xi$
algebraisch durch Radikale zu bestimmen.}

Wenn wir aber nicht nur eine, sondern sämtliche Wurzeln $\xi$ der
Transformationsgleichung adjungieren, oder, was auf dasselbe hinauskommt, die drei zu
$R_{1,0}$, $R_{0,1}$, $R_{1,1}$ gehörigen, so reduziert sich die Gruppe der
Teilungsgleichung auf $\mathfrak{A}_0$, und wenn wir noch $n$te Einheitswurzeln
adjungieren, auf $\mathfrak{B}_0$. Die Gruppe $\mathfrak{B}_0$ ist eine \gesperrt{Abelsche} vom
Grade $2^k$, die nur solche Elemente enthält, deren Grad $= 2$ ist. Infolgedessen ist die
\gesperrt{Teilungsgleichung durch Quadratwurzeln lösbar.}

Um die Form dieser Lösung zu finden, setze man
\[
n = p^\pi p'^{\pi'} \ldots,
\]
worin $p, p', \ldots$ voneinander verschiedene Primzahlen, $\pi, \pi', \ldots$
positive Exponenten sind. Man bestimme die Zahlen $c, c', \ldots$ aus den Kongruenzen
\[
\begin{aligned}
c &\equiv -1\ (\mathrm{mod}\ p^\pi), & c' &\equiv -1\ (\mathrm{mod}\ p'^{\pi'}) \ldots \\
c &\equiv +1\ (\mathrm{mod}\ np^{-\pi}), & c' &\equiv +1\ (\mathrm{mod}\ np'^{-\pi'}) \ldots
\end{aligned}
\]
dann erhält man jede Lösung $\alpha$ der Kongruenz
\[
\alpha^2 \equiv 1\ (\mathrm{mod}\ n),
\]
und jede nur einmal, in der Form
\weq{12}{\alpha \equiv c^\varepsilon c'^{\varepsilon'} \ldots \ (\mathrm{mod}\ n),}
wenn $\varepsilon, \varepsilon', \ldots$ die Werte 0, 1 annehmen.

Ist nun $\Omega$ irgend einer der Werte
\[
\frac{4\mu K + 4\mu' i K'}{n},
\]
%CONT>>
